% Unit 123 terjemahan Bahasa Indonesia dari chapter9.tex baris 580--817.
% Sumber: Methods of Algebra, Volume 2, commit
% 9a5803ff77dd3257484cb177f851a73770a59dd3 (CC BY 4.0).
% stable-unit-id: o014.aljabr2.chapter9.trace-and-dimension
% Cakupan: Bagian 9.2 lengkap, Dualitas: Jejak dan Dimensi, hingga tepat
% sebelum Bagian 9.3.
% Perbaikan matematis-redaksional: rumus penutup pada baris sumber 816
% diberi tanda kurung sehingga terbaca benar sebagai dim(X tensor Y),
% bukan (dim X) tensor Y.

% segment-id: o014.aljabr2.chapter9.trace-and-dimension.h001
\section{Dualitas: Jejak dan Dimensi}\label{sec:symm-monoidal-dual}
\hypertarget{o014.aljabr2.chapter9.trace-and-dimension}{ }
% segment-id: o014.aljabr2.chapter9.trace-and-dimension.p001
Dalam \S\ref{sec:monoidal-dual} kita telah memperkenalkan beberapa konsep
yang berkaitan dengan dualitas. Semuanya memiliki versi kiri dan kanan,
bergantung pada posisi objek di dalam $\otimes$. Dalam banyak penerapan,
kategori monoidal yang ditinjau bersifat simetris sehingga kedua versi itu
tidak perlu dibedakan. Kita mulai dengan kasus kategori monoidal beranyam.
Seperti biasa, struktur kepangnya ditulis dalam bentuk
% segment-id: o014.aljabr2.chapter9.trace-and-dimension.q001
\[ c(X,Y): X \otimes Y \rightiso Y \otimes X \]
Struktur kepang tersebut simetris tepat ketika
$c(X,Y)^{-1}=c(Y,X)$. Diagram kepang dapat digunakan bersama teknik diagram
tali yang diperkenalkan dalam \S\ref{sec:monoidal-dual}.

% segment-id: o014.aljabr2.chapter9.trace-and-dimension.l001
\begin{lemma}\label{prop:braided-dual}
	Misalkan $\mathcal{C}$ kategori monoidal beranyam. Jika $X^*$ adalah dual
	kanan dari $X\in\Obj(\mathcal{C})$, maka $X^*$ juga merupakan dual kiri
	dari $X$. Lebih tepatnya, definisikan
	% segment-id: o014.aljabr2.chapter9.trace-and-dimension.q002
	\begin{align*}
		\mathrm{ev}' & := \mathrm{ev} \circ c(X, X^*)^{-1}: X^* \otimes X \to \munit, \\
		\mathrm{coev}' & := c(X^*, X) \circ \mathrm{coev}: \munit \to X \otimes X^* ,
	\end{align*}
	maka $(X^*,X,\mathrm{ev}',\mathrm{coev}')$ adalah data dualitas. Kasus
dual kiri ${}^*X$ serupa. Jika $\mathcal{C}$ kategori monoidal simetris,
berlaku pula $\mathrm{ev}''=\mathrm{ev}$ dan
$\mathrm{coev}''=\mathrm{coev}$.
\end{lemma}
% segment-id: o014.aljabr2.chapter9.trace-and-dimension.pf001
\begin{proof}
	Ini merupakan penerapan aksioma struktur kepang. Sebagai contoh,
	$(\identity\otimes\mathrm{ev}')(\mathrm{coev}'\otimes\identity)
	=\identity$ direduksi menjadi komutativitas diagram berikut:
	% segment-id: o014.aljabr2.chapter9.trace-and-dimension.g001
	\[\begin{tikzcd}
		X \otimes \munit \arrow[d, "{c(X, \munit)}"] \arrow[r, "{\identity \otimes \mathrm{coev}}"] & X \otimes X^* \otimes X \arrow[d, "{c(X, X^* \otimes X)}"] \arrow[rrr, "{\mathrm{ev} \otimes \identity}"] & & & \munit \otimes X \arrow[d, "{c(\munit, X)}"'] \\
		\munit \otimes X \arrow[r, "{\mathrm{coev} \otimes \identity}"' inner sep=0.6em] & X^* \otimes X \otimes X \arrow[r, "{c(X^*, X) \otimes \identity}"' inner sep=0.6em] & X \otimes X^* \otimes X \arrow[r, "{\identity \otimes c(X, X^*)^{-1}}"' inner sep=0.6em] & X \otimes X \otimes X^* \arrow[r, "{\identity \otimes \mathrm{ev}}"' inner sep=0.6em] & X \otimes \munit
	\end{tikzcd}\]
	Persegi kecil komutatif karena funktorialitas struktur kepang, sedangkan
persegi besar di sebelah kanan merupakan argumen diagram kepang yang
standar; bandingkan \S\ref{sec:alg-in-monoidal-cat}.
\end{proof}

% segment-id: o014.aljabr2.chapter9.trace-and-dimension.p002
Jadi, sebuah objek $X$ dalam kategori monoidal beranyam memiliki dual kiri
jika dan hanya jika ia memiliki dual kanan. Dalam kasus kategori monoidal
simetris, kita dapat mengidentifikasi kedua versi itu dan menulis objek
serta morfisme dual secara seragam sebagai $X^*$ dan $f^*$.

% segment-id: o014.aljabr2.chapter9.trace-and-dimension.p003
Contoh yang paling khas tentu saja ialah kategori modul di atas gelanggang
komutatif. Dengan hasil kali tensor modul, kategori ini merupakan kategori
monoidal simetris, dan dualitas di dalamnya memiliki karakterisasi sederhana.

% segment-id: o014.aljabr2.chapter9.trace-and-dimension.pr001
\begin{proposition}\label{prop:dualizable-module}
	Misalkan $R$ gelanggang komutatif. Jadikan $R\dcate{Mod}$ kategori
	monoidal simetris melalui $\otimes_R$. Suatu $R$-modul $M$ memiliki
	dual $M^*$ jika dan hanya jika $M$ merupakan modul projektif yang
	dibangkitkan secara hingga. Dalam hal ini dapat diambil
	$M^*:=\Hom_R(M,R)$, sedangkan $\mathrm{ev}_M$ dan $\mathrm{coev}_M$
	diberikan seperti dalam Definisi~\ref{def:bimodule-ev-coev} (ambil
	$A=B=\Bbbk=R$ dan $P=M$; $M^*$ di sini dinotasikan sebagai $M^\vee$
	di sana), mungkin setelah menukar urutan faktor tensor.
\end{proposition}
% segment-id: o014.aljabr2.chapter9.trace-and-dimension.pf002
\begin{proof}
	Arah ``jika'' langsung karena pilihan
	$(M^*,\mathrm{ev}_M,\mathrm{coev}_M)$ bersifat eksplisit. Dalam kasus
	khusus $M=R$, kita dapat mengidentifikasi $M^*$ dengan $R$ dan memperoleh
	$\mathrm{ev}_M(x\otimes y)=xy$ serta
	$\mathrm{coev}_M(1)=1\otimes1$; semua identitas yang diperlukan lalu
	terpenuhi secara langsung.

	Konstruksi $(M^*,\mathrm{ev}_M,\mathrm{coev}_M)$ jelas kompatibel dengan
jumlah langsung, sehingga diperoleh kasus $M=R^{\oplus n}$. Setiap modul
projektif yang dibangkitkan secara hingga $M$ dapat direalisasikan sebagai
suku langsung dari suatu $R^{\oplus n}$; jadi $M$ memiliki dual.

	Sekarang tinjau arah ``hanya jika''. Misalkan $R$-modul $M$ memiliki
dual $M^*$. Berdasarkan Proposisi~\ref{prop:dual-internal-Hom} (ii), untuk
setiap $R$-modul $M'$ terdapat isomorfisme $R$-modul
	% segment-id: o014.aljabr2.chapter9.trace-and-dimension.g002
	\[\begin{tikzcd}[row sep=tiny]
		\Hom_R(M, M') & \Hom_R(R, M^* \dotimes{R} M') \arrow[l, "\sim"'] \arrow[r, "\sim"] & M^* \dotimes{R} M' \\
		(\mathrm{ev}_M \otimes \identity_{M'}) (\identity_M \otimes \varphi) & \varphi \arrow[mapsto, l] \arrow[mapsto, r] & \varphi(1).
	\end{tikzcd}\]

	Dengan mengambil $M'=R$, kita memperoleh
	$\Hom_R(M,R)\simeq\Hom_R(R,M^*)\simeq M^*$. Isomorfisme itu mengirim
$\lambda\in M^*$ ke homomorfisme berikut:
	% segment-id: o014.aljabr2.chapter9.trace-and-dimension.g003
	\[\begin{tikzcd}[row sep=tiny]
		M \arrow[r, "\sim"] & M \dotimes{R} R \arrow[r, "{\identity_M \otimes \lambda}"] & M \dotimes{R} M^* \arrow[r, "{\mathrm{ev}_M}"] & R \\ 
		m \arrow[mapsto, r] & m \otimes 1 \arrow[mapsto, r] & m \otimes \lambda \arrow[mapsto, r] & \mathrm{ev}_M(m \otimes \lambda).
	\end{tikzcd}\]
	Dengan demikian, $M^*$ dapat diidentifikasi dengan $\Hom_R(M,R)$ dan
$\mathrm{ev}_M$ dengan evaluasi.

	Tinjau identitas
	$(\mathrm{ev}_M\otimes\identity_M)
	(\identity_M\otimes\mathrm{coev}_M)=\identity_M$. Tuliskan
	$\mathrm{coev}_M(1)=\sum_{i=1}^n\lambda_i\otimes m_i$. Identitas itu
	setara dengan $\sum_i\lambda_i(m)m_i=m$ untuk setiap $m$, yakni
$\identity_M$ terfaktorkan sebagai
	% segment-id: o014.aljabr2.chapter9.trace-and-dimension.q003
	\[ M \xrightarrow{(\lambda_1, \ldots, \lambda_n)} R^{\oplus n} \xrightarrow{(m_1, \ldots, m_n)} M, \]
	Hal ini menunjukkan bahwa $M$ merupakan suku langsung dari
$R^{\oplus n}$.
\end{proof}

% segment-id: o014.aljabr2.chapter9.trace-and-dimension.d001
\begin{definition}\label{def:rigid-cat-braided}
	\index{kategori monoidal!kaku (rigid)}
	Untuk kategori monoidal beranyam, kekakuan kiri dan kekakuan kanan dalam
	Definisi~\ref{def:rigid-cat} saling ekuivalen. Kategori monoidal beranyam
	yang memenuhi salah satunya disebut \emph{kategori kaku (rigid)}.
\end{definition}

% segment-id: o014.aljabr2.chapter9.trace-and-dimension.p004
Setelah dilengkapi dengan struktur kategori abelian, objek unit dari kategori
monoidal beranyam yang kaku memperlihatkan sifat khusus. Pertama-tama kita
buktikan bahwa $\otimes$ otomatis aditif pada setiap variabel.

% segment-id: o014.aljabr2.chapter9.trace-and-dimension.pr002
\begin{proposition}\label{prop:rigid-additive}
	Misalkan $\mathcal{C}$ kategori monoidal beranyam yang kaku. Jika
	$\mathcal{C}$ juga kategori aditif, maka $\otimes$ adalah bifunktor aditif.
\end{proposition}
% segment-id: o014.aljabr2.chapter9.trace-and-dimension.pf003
\begin{proof}
	Pilih $X\in\Obj(\mathcal{C})$. Proposisi~\ref{prop:dual-internal-Hom}
	(iii) menyiratkan bahwa funktor $X\otimes(\cdot)$ memiliki adjoin kanan
$X^*\otimes(\cdot)$. Karena itu, Korolari~\ref{prop:automatic-additivity}
(v) menjamin bahwa $X\otimes(\cdot)$ aditif. Berdasarkan struktur kepang,
hal yang sama berlaku bagi $(\cdot)\otimes X$.
\end{proof}

% segment-id: o014.aljabr2.chapter9.trace-and-dimension.pr003
\begin{proposition}\label{prop:rigid-unit-simple}
	Misalkan $\mathcal{C}$ kategori monoidal beranyam yang kaku dan sekaligus
kategori abelian. Jika $\iota:U\hookrightarrow\munit$ adalah subobjek, maka
$\munit=U\oplus\Ker(\iota^*)$. Akibatnya:
	\begin{enumerate}[(i)]
		\item objek $\munit$ terbelah (Definisi~\ref{def:semisimple});
		\item jika $\End_{\mathcal{C}}(\munit)$ sebuah medan, maka $\munit$
		adalah objek sederhana.
	\end{enumerate}
\end{proposition}
% segment-id: o014.aljabr2.chapter9.trace-and-dimension.pf004
\begin{proof}
	Definisikan $V:=\Coker(\iota)$. Telah diketahui bahwa $\otimes$ merupakan
funktor eksak aditif pada setiap variabel (Korolari~\ref{prop:rigid-exact}
dan Proposisi~\ref{prop:rigid-additive}). Karena itu, kita memperoleh diagram
komutatif berikut, yang baris-baris pada bagian garis penuhnya eksak:
	% segment-id: o014.aljabr2.chapter9.trace-and-dimension.g004
	\[\begin{tikzcd}
		0 \arrow[r] & U \arrow[r, "\iota"] & \munit \arrow[r] & V \arrow[r] & 0 \\
		0 \arrow[r] & U \otimes U \arrow[hookrightarrow, u, "{\iota \otimes \identity}"] \arrow[r, "{\identity \otimes \iota}"'] & U \arrow[hookrightarrow, u, "\iota"] \arrow[r] \arrow[dashed, ru] & U \otimes V \arrow[hookrightarrow, u, "{\iota \otimes \identity}"'] \arrow[r] & 0
	\end{tikzcd}\]
	Komposit panah putus-putus adalah $0$. Oleh karena itu,
$U\otimes V=0$ dan $U\otimes U\simeq U$.

	Untuk setiap objek $T$, dengan cara yang sama diperoleh monomorfisme
$\iota\otimes\identity_T:U\otimes T\hookrightarrow T$. Jadi
$U\otimes T=0$ jika dan hanya jika
$\iota\otimes\identity_T=0$. Pernyataan terakhir ini pada gilirannya
	ekuivalen dengan lenyapnya morfisme $T\to U^*\otimes T$ yang bersesuaian
di bawah Proposisi~\ref{prop:dual-internal-Hom} (iii). Tidak sulit
menunjukkan bahwa morfisme terakhir tersebut tepatnya
$\iota^*\otimes\identity_T$ (Latihan bab ini). Dengan demikian, untuk
setiap objek $X$, subobjek terbesar $T\subset X$ yang memenuhi
$U\otimes T=0$ sama dengan subobjek terbesar yang membuat
$T\to U^*\otimes T\hookrightarrow U^*\otimes X$ bernilai nol, dan ini
juga sama dengan
	% segment-id: o014.aljabr2.chapter9.trace-and-dimension.q004
	\[ \Ker\left[ \iota^* \otimes \identity_X: X \to U^* \otimes X \right] \simeq \Ker(\iota^*) \otimes X. \]
	\begin{itemize}
		\item Terapkan hal ini pada $X=V$ dan gunakan $U\otimes V=0$;
		kita memperoleh $\Ker(\iota^*)\otimes V\simeq V$.
		\item Terapkan hal ini pada $X=U$. Karena setiap subobjek $T$ dari
		$U$ juga merupakan subobjek dari $\munit$, berlaku
		$U\otimes T\supset T\otimes T\simeq T$. Akibatnya,
		$\Ker(\iota^*)\otimes U=0$.
	\end{itemize}

	Substitusikan kedua hasil ini ke dalam barisan eksak pendek
	% segment-id: o014.aljabr2.chapter9.trace-and-dimension.q005
	\[ 0 \to \Ker(\iota^*) \otimes U \to \Ker(\iota^*) \to \Ker(\iota^*) \otimes V \to 0, \]
	maka langsung diperoleh
$\munit\supset\Ker(\iota^*)\rightiso V$. Hal ini membelah barisan
$0\to U\to\munit\to V\to0$. Karena $U$ sembarang subobjek, inilah tepatnya
definisi bahwa $\munit$ terbelah.

	Terakhir, Korolari~\ref{prop:indecomposable-criterion} menyiratkan bahwa
jika $\End_{\mathcal{C}}(\munit)$ merupakan medan, maka $\munit$ tak
terurai; jadi $\munit$ sederhana.
\end{proof}

% segment-id: o014.aljabr2.chapter9.trace-and-dimension.p005
Dalam setiap kategori monoidal $\mathcal{C}$, $\End(\munit)$ merupakan
monoid dan mudah dibuktikan komutatif; lihat \cite[Latihan Bab 3]{Li1}.
Kendala unit $X\simeq\munit\otimes X$ membuat setiap unsur
$\End(\munit)$ bertindak pada setiap $X$ melalui endomorfisme. Endomorfisme
ini menentukan homomorfisme monoid
$\End(\munit)\to Z(\mathcal{C})$, dengan $Z(\mathcal{C})$ menyatakan pusat
kategori $\mathcal{C}$. Dari sifat umum pusat, komposisi morfisme bersifat
bilinear terhadap aksi $\End(\munit)$:
% segment-id: o014.aljabr2.chapter9.trace-and-dimension.q006
\[ (z f) g = z(f g) = f (z g), \quad X \xrightarrow{g} Y \xrightarrow{f} Z, \quad z \in \End(\munit). \]

% segment-id: o014.aljabr2.chapter9.trace-and-dimension.p006
Selain itu, dari aksioma $\otimes$ (lihat \cite[\S 3.1]{Li1}) mudah
diturunkan bahwa $\otimes$ juga bilinear:
% segment-id: o014.aljabr2.chapter9.trace-and-dimension.q007
\[ (zf_1) \otimes f_2 = z(f_1 \otimes f_2) = f_1 \otimes (z f_2), \quad f_i \in \Hom(X_i, Y_i), \quad z \in \End(\munit). \]

% segment-id: o014.aljabr2.chapter9.trace-and-dimension.p007
Dengan demikian, $\End(\munit)$ dapat dipandang sebagai semacam
``koefisien'' bagi kategori monoidal $\mathcal{C}$.

% segment-id: o014.aljabr2.chapter9.trace-and-dimension.d002
\begin{definition}
	\index{jejak (trace)}
	Misalkan $f\in\End(X)$ endomorfisme dalam kategori monoidal beranyam
$\mathcal{C}$ dan $X$ memiliki dual. Definisikan jejak kiri dari $f$
(disingkat \emph{jejak}) dengan
	% segment-id: o014.aljabr2.chapter9.trace-and-dimension.q008
	\[ \Tr(f) = \Tr_{\text{kiri}}(f) := \mathrm{ev}_X \left( f \otimes \identity_{X^*} \right) \mathrm{coev}'_X \; \in \End(\munit), \]
	dan definisikan jejak kanannya dengan
	% segment-id: o014.aljabr2.chapter9.trace-and-dimension.q009
	\[ \Tr_{\text{kanan}}(f) := \mathrm{ev}'_X \left( \identity_{X^*} \otimes f \right) \mathrm{coev}_X, \]
	dengan $\mathrm{coev}'_X$ dan $\mathrm{ev}'_X$ seperti dalam
Lema~\ref{prop:braided-dual}. Berdasarkan Definisi--Proposisi
\ref{def:dual-uniqueness}, kedua jejak tidak bergantung pada pilihan data
dualitas.
\end{definition}

% segment-id: o014.aljabr2.chapter9.trace-and-dimension.d003
\begin{definition}
	\index{dimensi (dimension)}
	Misalkan $X$ objek dalam kategori monoidal beranyam $\mathcal{C}$ dan
$X$ memiliki dual. Dimensi kirinya (disingkat \emph{dimensi})
didefinisikan sebagai
	% segment-id: o014.aljabr2.chapter9.trace-and-dimension.q010
	\[ \dim X = \dim_{\text{kiri}} X := \Tr(\identity_X) \; \in \End(\munit), \]
	dan dimensi kanannya didefinisikan sebagai
$\dim_{\text{kanan}}X:=\Tr_{\text{kanan}}(\identity_X)$.
\end{definition}

% segment-id: o014.aljabr2.chapter9.trace-and-dimension.g005
Jejak kiri dan kanan masing-masing digambarkan sebagai berikut:
\[\Tr_{\text{kiri}}(f) = \begin{tikzpicture}[baseline=(M), bend angle=70]
	\node (A) {$X$};
	\node (B) [right=of A] {$X^*$};
	\node (C) [below=of A] {$X$};
	\node (D) [below=of B] {$X^*$};
	
	\path (A) edge[bend left] node[auto] {$\mathrm{coev}'_X$} (B);
	\path (B) edge (D);
	\path (A) edge node[midway, name=M, draw, fill=white] {$f$} (C);
	\path (D) edge[bend left] node[auto, swap] {$\mathrm{ev}_X$} (C);
\end{tikzpicture}
\xlongequal{\text{disingkat}}\;
\begin{tikzpicture}[baseline=(M)]
	\coordinate (A) at (0, 0);
	\coordinate (B) at (1, 0);
	\coordinate (C) at (0, -1);
	\coordinate (D) at (1, -1);
	\draw (A) arc[start angle=180, end angle=0, radius=0.5] -- (D) arc[start angle=360, end angle=180, radius=0.5] -- (A);
	\node[draw, fill=white] (M) at ($(A)!.5!(C)$) {$f$};
\end{tikzpicture} \qquad
\Tr_{\text{kanan}}(f) \xlongequal{\text{disingkat}}\; \begin{tikzpicture}[baseline=(M)]
	\coordinate (A) at (0, 0);
	\coordinate (B) at (1, 0);
	\coordinate (C) at (0, -1);
	\coordinate (D) at (1, -1);
	\draw (A) arc[start angle=175, end angle=0, radius=0.5] -- (D) arc[start angle=360, end angle=180, radius=0.5] -- (A);
	\node[draw, fill=white] (M) at ($(B)!.5!(D)$) {$f$};
\end{tikzpicture}\]

% segment-id: o014.aljabr2.chapter9.trace-and-dimension.p008
Jika $\mathcal{C}$ kategori monoidal simetris, versi kiri dan kanan dari
jejak maupun dimensi selalu sama. Inilah situasi utama yang akan digunakan
kelak.

% segment-id: o014.aljabr2.chapter9.trace-and-dimension.p009
Jika peran kategori $\mathcal{C}$ perlu ditonjolkan, kita menggunakan
notasi seperti $\Tr_{\mathcal{C}}$ dan $\dim_{\mathcal{C}}$. Dalam sebagian
literatur, keduanya juga disebut jejak kuantum dan dimensi kuantum.
Contoh~\ref{eg:Vect-trace} akan menjelaskan hubungannya dengan versi klasik.

% segment-id: o014.aljabr2.chapter9.trace-and-dimension.pr004
\begin{proposition}
	Misalkan $F:\mathcal{C}\to\mathcal{D}$ funktor monoidal antara kategori
	monoidal beranyam yang mempertahankan struktur kepang. Untuk setiap
	endomorfisme $f\in\End(X)$ di $\mathcal{C}$, dengan $X$ memiliki dual,
	berlaku
$F\left(\Tr_{\mathcal{C}}(f)\right)=\Tr_{\mathcal{D}}(Ff)$; demikian pula
untuk jejak kanan. Khususnya,
$F\left(\dim_{\mathcal{C}}X\right)=\dim_{\mathcal{D}}(FX)$.
\end{proposition}
% segment-id: o014.aljabr2.chapter9.trace-and-dimension.pf005
\begin{proof}
	Buktinya sama persis dengan Proposisi~\ref{prop:monoidal-functor-dual}.
\end{proof}

% segment-id: o014.aljabr2.chapter9.trace-and-dimension.p010
Jejak yang didefinisikan di atas memiliki sifat-sifat yang serupa dengan
jejak pemetaan linear.

% segment-id: o014.aljabr2.chapter9.trace-and-dimension.pr005
\begin{proposition}\label{prop:trace}
	Misalkan $\mathcal{C}$ kategori monoidal beranyam. Dalam pernyataan
berikut, setiap kali dibahas endomorfisme $f\in\End(X)$, selalu diasumsikan
bahwa $X$ memiliki dual.
	\begin{enumerate}[(i)]
		\item Berlaku
		$\Tr_{\text{kiri}}(f)=\Tr_{\text{kanan}}({}^*f)$ dan
		$\Tr_{\text{kanan}}(f)=\Tr_{\text{kiri}}(f^*)$.
		\item Misalkan $X$ dan $Y$ keduanya memiliki dual. Untuk setiap
		$\begin{tikzcd} X \arrow[shift left, r, "f"] & Y \arrow[shift left, l, "g"] \end{tikzcd}$,
		berlaku
		$\Tr_{\text{kiri}}(gf)=\Tr_{\text{kiri}}(({}^{**}f)g)$ dan
		$\Tr_{\text{kanan}}(gf)=\Tr_{\text{kanan}}((f^{**})g)$.
		Perhatikan bahwa jika $\mathcal{C}$ kategori monoidal simetris, maka
		${}^{**}f=f=f^{**}$.
		\item Untuk setiap $f\in\End(X)$ dan $g\in\End(Y)$, berlaku
		$\Tr(f\otimes g)=\Tr(f)\Tr(g)$.
		\item Jika $a\in\End(\munit)$, maka selalu berlaku
		$\Tr(af)=a\Tr(f)$.
		\item Misalkan $\mathcal{C}$ kategori-$\cate{Ab}$ dan $\otimes$
		aditif pada variabel pertama (atau kedua). Maka untuk setiap
		$f,g\in\End(X)$ berlaku
		$\Tr_{\text{kiri}}(f+g)=\Tr_{\text{kiri}}(f)+\Tr_{\text{kiri}}(g)$
		(atau
		$\Tr_{\text{kanan}}(f+g)=\Tr_{\text{kanan}}(f)+\Tr_{\text{kanan}}(g)$).
	\end{enumerate}
	Dalam pernyataan (iii) dan (iv), $\Tr$ dapat diartikan sebagai jejak kiri
maupun jejak kanan.
\end{proposition}
% segment-id: o014.aljabr2.chapter9.trace-and-dimension.pf006
\begin{proof}
	Seperti biasa, pembuktian terutama mengandalkan gambar. Karena kasus kiri
dan kanan simetris, cukup membahas jejak kiri. Pertama, berdasarkan
diagram jejak, menutup bagian bawah setiap suku dalam
\eqref{eqn:f-dual-coev} langsung memberikan (i). Dengan cara serupa, (ii)
digambarkan melalui \eqref{eqn:f-dual-ev} sebagai berikut.
	% segment-id: o014.aljabr2.chapter9.trace-and-dimension.g006
	\[\begin{tikzpicture}[baseline=(O)]
		\coordinate (A) at (0, 0);
		\coordinate (B) at (1, 0);
		\coordinate (C) at (0, -1);
		\coordinate (D) at (1, -1);
		\coordinate (O) at (0, -0.5);
		\draw (A) arc[start angle=180, end angle=0, radius=0.5] -- (D) arc[start angle=360, end angle=180, radius=0.5] -- (A);
		\node[draw, fill=white] (M) at ($(A)!.2!(C)$) {$f$};
		\node[draw, fill=white] (N) at ($(A)!.8!(C)$) {$g$};
	\end{tikzpicture} = \begin{tikzpicture}[baseline=(O)]
		\coordinate (A) at (0, 0);
		\coordinate (B) at (1, 0);
		\coordinate (C) at (0, -1);
		\coordinate (D) at (1, -1);
		\coordinate (O) at (0, -0.5);
		\draw (A) arc[start angle=180, end angle=0, radius=0.5] -- (D) arc[start angle=360, end angle=180, radius=0.5] -- (A);
		\node[draw, fill=white] (M) at ($(B)!.5!(D)$) {${}^* f$};
		\node[draw, fill=white] (N) at ($(A)!.5!(C)$) {$g$};
	\end{tikzpicture} = \begin{tikzpicture}[baseline=(O)]
		\coordinate (A) at (0, 0);
		\coordinate (B) at (1, 0);
		\coordinate (C) at (0, -1);
		\coordinate (D) at (1, -1);
		\coordinate (O) at (0, -0.5);
		\draw (A) arc[start angle=180, end angle=0, radius=0.5] -- (D) arc[start angle=360, end angle=180, radius=0.5] -- (A);
		\node[draw, fill=white] (M) at ($(A)!.2!(C)$) {$g$};
		\node[draw, fill=white] (N) at ($(A)!.8!(C)$) {${}^{**} f$};
	\end{tikzpicture}\]

	Untuk (iii), intinya ialah identitas intuitif
	% segment-id: o014.aljabr2.chapter9.trace-and-dimension.g007
	\[\begin{tikzpicture}[baseline=(O)]
		\coordinate (O) at (0, 0);
		\draw (O) circle[radius=1];
		\draw (O) circle[radius=0.4];
		\node[draw, fill=white] at ($(O) + (180:1)$) {$f$};
		\node[draw, fill=white] at ($(O) + (180:0.4)$) {$g$};
	\end{tikzpicture} \; = \;
	\begin{tikzpicture}[baseline=(A)]
		\coordinate (A) at (0, 0);
		\coordinate (B) at (1.2, 0);
		\draw (A) circle[radius=0.4];
		\draw (B) circle[radius=0.4];
		\node[draw, fill=white] at ($(A) + (180:0.4)$) {$f$};
		\node[draw, fill=white] at ($(B) + (180:0.4)$) {$g$};
	\end{tikzpicture} \; = \;
	\begin{tikzpicture}[baseline=(O)]
		\coordinate (A) at (0, 1);
		\coordinate (B) at (0, 0);
		\coordinate (O) at ($(A)!.5!(B)$);
		\draw (A) circle[radius=0.4];
		\draw (B) circle[radius=0.4];
		\node[draw, fill=white] at ($(A) + (180:0.4)$) {$f$};
		\node[draw, fill=white] at ($(B) + (180:0.4)$) {$g$};
	\end{tikzpicture}\]
	Bukti formalnya didasarkan pada berbagai sifat kealamian objek unit.

	Pernyataan (iv) mengikuti dari bilinearitas komposisi morfisme dan
$\otimes$ terhadap $\End(\munit)$. Pernyataan (v) mengikuti langsung dari
definisi tertulis jejak kiri dan kanan, tanpa memerlukan gambar.
\end{proof}

% segment-id: o014.aljabr2.chapter9.trace-and-dimension.p011
Proposisi~\ref{prop:trace} (iv) dan (v) dapat dipahami sebagai sifat
linearitas jejak. Sebagai kasus khusus (iii), berlaku pula
$\dim(X\otimes Y)=\dim X\,\dim Y$.
